\documentclass[10pt,conference]{IEEEtran}
\usepackage{amsmath,amssymb,booktabs,array}
\usepackage[hidelinks]{hyperref}

\title{A Parametric Chisel Architecture for Fixed-Graph Quantum LDPC Min-Sum Decoders}
\author{Anonymous Authors}

\begin{document}
\maketitle

\begin{abstract}
This manuscript records the design and validation of \texttt{chipsLDPC}, a
small Chisel/CIRCT research generator for field-programmable gate array
implementations of quantum low-density parity-check decoders.  The current
artifact implements shared check-node and variable-node primitives for the
syndrome-based min-sum architecture of Valls et al. and the memory-augmented
architecture of Maurer et al.  The implementation deliberately separates
static graph elaboration from runtime message processing, represents edge
messages using distinct sign and magnitude fields, and provides executable
tests for ties, zero, fixed-width arithmetic, saturation, and cycle latency.
An immutable pure Scala graph model validates sparse parity-check adjacency and
assigns deterministic edge identifiers before hardware elaboration.
Independent binary Gauss--Jordan processing elements provide a small
post-processing interoperability example.
This document describes the implemented milestone; performance and full-code
decoder results are reserved for future revisions.
\end{abstract}

\section{Research Objective}
Quantum LDPC decoders must combine algorithmic experimentation with hardware
structures whose resource usage and timing remain transparent.  The objective
of \texttt{chipsLDPC} is therefore not to provide a large decoding framework.
It is to provide a small family of composable generators from which a fixed
decoder graph can be elaborated, inspected as FIRRTL-dialect MLIR, emitted as
synthesizable SystemVerilog, simulated with Verilator, and eventually evaluated
on an FPGA.

The current milestone asks a narrow question: can the CNU and VNU datapaths in
Refs.~\cite{valls2021,maurer2025} be represented by one common implementation
with explicit, testable policy differences?  Controller construction, decoder
termination, sliding windows, syndrome streaming, and FPGA implementation are
outside the present scope.

\section{Algorithmic Background}
Let $H\in\{0,1\}^{m\times n}$ be a parity-check matrix.  Check $i$ is adjacent
to variables $\mathcal{N}(i)$ and variable $j$ is adjacent to checks
$\mathcal{M}(j)$.  The incoming variable-to-check message is denoted
$R_{i,j}$ and the outgoing check-to-variable message is denoted
$\sigma_{i,j}$.  A syndrome bit $S_i$ modifies the check parity.

For syndrome-based min-sum, the check update is
\begin{equation}
 \sigma_{i,j} =
 \left(
 S_i \oplus
 \bigoplus_{j'\in\mathcal{N}(i)\setminus j}
 \operatorname{sgn}(R_{i,j'})
 \right)
 \alpha
 \min_{j'\in\mathcal{N}(i)\setminus j}|R_{i,j'}|.
 \label{eq:cnu}
\end{equation}
The corresponding variable update is
\begin{equation}
 R_{i,j}=\Lambda_j+
 \sum_{i'\in\mathcal{M}(j)\setminus i}\sigma_{i',j},
 \label{eq:vnu}
\end{equation}
with marginal
\begin{equation}
 M_j=\Lambda_j+\sum_{i'\in\mathcal{M}(j)}\sigma_{i',j}.
 \label{eq:marginal}
\end{equation}
The hard decision is
\begin{equation}
 \operatorname{HD}(M_j)=
 \begin{cases}
  1,&M_j<0,\\
  0,&M_j\geq 0.
 \end{cases}
 \label{eq:hard-decision}
\end{equation}
This strict-negative convention makes zero a no-error decision and agrees with
the canonical sign--magnitude representation, in which zero always has sign
zero.  The decision is taken from the full working-width marginal before
saturation.

The memory-augmented VNU replaces the static prior by
\begin{equation}
 \Lambda_j(t)=(1-\gamma_j)\Lambda_j(0)+\gamma_j M_j(t-1).
 \label{eq:memory}
\end{equation}
Defining $\beta_j=1-\gamma_j$ gives the implemented form
\begin{equation}
 \Lambda_j(t)=\beta_j\Lambda_j(0)+M_j(t-1)-\beta_jM_j(t-1).
 \label{eq:beta}
\end{equation}
This form permits one unsigned fixed-point coefficient to multiply two signed
operands.  Coefficient generation is intentionally external to the VNU.

\section{Design Principles}
\subsection{Elaboration and runtime data}
The parity-check matrix and adjacency lists are Scala values used during
elaboration.  A future full decoder will construct one CNU per row, one VNU per
column, and fixed wires per nonzero of $H$.  The matrix is not a runtime port.
Syndromes, messages, iteration controls, and memory coefficients are runtime
hardware data.

\subsection{Static Tanner-graph contract}
For $H\in\{0,1\}^{m\times n}$, the implemented \texttt{TannerGraph} accepts
$n$ and \texttt{rowOnes}, where \texttt{rowOnes(i)} contains the variable
indices in $\mathcal{N}(i)$.  Each row is sorted, and global edge identifiers
are then assigned in increasing check order and increasing variable order.
Thus a physical message port has one deterministic identity independent of the
input order within a row.

The column adjacency \texttt{colOnes}, check-to-edge mapping, and
variable-to-edge mapping are derived from that single normalized source.  The
constructor rejects nonpositive dimensions, empty check or variable adjacency,
duplicates, and indices outside $[0,n)$.  It does not impose a min-sum-specific
minimum check degree: algorithm constraints belong to the later hardware
elaborator.  The graph contains no Chisel types and produces no registers,
ports, memories, or runtime crossbar.

\texttt{TannerNodeGraphs} derives one rooted view per check and per variable.
Each view retains the global edge identifier, while its vector position defines
the corresponding local CNU or VNU port.  A check view records the neighboring
variable for each edge; a variable view records the neighboring check.  The
views therefore duplicate only immutable Scala metadata and cannot introduce a
second edge numbering or any runtime routing structure.

The concrete validation example uses one CSS sector of the Steane
$[[7,1,3]]$ code~\cite{steane1996}.  With zero-based variable indices, the
chosen Hamming parity-check convention is
\begin{equation}
 H_{\mathrm{Steane}} =
 \begin{bmatrix}
 0&0&0&1&1&1&1\\
 0&1&1&0&0&1&1\\
 1&0&1&0&1&0&1
 \end{bmatrix},
\end{equation}
so \texttt{rowOnes} is
$[[3,4,5,6],[1,2,5,6],[0,2,4,6]]$.  The code has identical Hamming-code
blocks $H_X=H_Z$, so the same Tanner topology applies independently to either
CSS decoding sector.

\subsection{Composition rather than module inheritance}
Immutable Scala case classes describe degrees, widths, and policies.  Chisel
bundles describe interfaces.  Pure generator functions implement small
combinational structures, and modules provide meaningful hierarchy or state.
Paper variants are elaboration-time policies around a shared datapath rather
than subclasses containing duplicated hardware.

\subsection{Inspectability}
Combinational implementations are \texttt{RawModule}s and therefore introduce
no implicit clock or reset.  Registered wrappers make the one-cycle CNU and VNU
phase boundaries explicit.  The generator emits both FIRRTL-dialect MLIR and
SystemVerilog so arithmetic widths, muxes, comparisons, and registers can be
inspected before synthesis.

\section{Numeric Contract}
Table~\ref{tab:contract} is normative for the current implementation.

\begin{table}[t]
\caption{Implemented fixed-width semantics.}
\label{tab:contract}
\centering
\begin{tabular}{@{}p{0.30\columnwidth}p{0.64\columnwidth}@{}}
\toprule
Property & Contract \\
\midrule
Edge representation & One sign bit and an unsigned magnitude. \\
Sign convention & Zero is nonnegative; one is negative. \\
Canonical zero & A zero magnitude always forces sign zero. \\
Minimum tie & Equal values occupy both minimum positions. \\
Selector & \texttt{useSecond} iff the incoming magnitude equals the first minimum. \\
Accumulator & Signed two's-complement with explicit guard width. \\
Decision tie & A zero marginal produces correction zero. \\
Saturation & Applied at named interfaces, not silently at tree levels. \\
\bottomrule
\end{tabular}
\end{table}

The edge magnitude width $w_m$ and stored accumulator width $w_a$ are explicit
configuration values.  The VNU evaluates its balanced sum at
\begin{equation}
 w_{\mathrm{work}} = w_a +
 \max\left(1,\left\lceil\log_2(d_v+1)\right\rceil\right),
\end{equation}
where $d_v$ is the variable degree.  The marginal is saturated into $w_a$,
while each extrinsic result is saturated into the unsigned $w_m$ magnitude.
Expanding or truncating arithmetic is selected explicitly in Chisel.
The locked Relay-oriented default is $(w_m,w_a)=(4,5)$, with four magnitude
bits, a separate message sign, and five-bit two's-complement stored marginals.

\section{Check-Node Unit}
\subsection{Exclusive sign}
For one check, the implementation first canonicalizes every input sign and
forms one global XOR with the syndrome.  The sign for edge $j$ is obtained by
XORing the global parity with the canonical sign of edge $j$.  This shares one
parity tree across all outputs.  This parity metadata remains independent of
the broadcast magnitude and is canonicalized after the VNU selects a minimum.

\subsection{Two-minimum tournament}
Every outgoing magnitude is either the smallest or second-smallest incoming
magnitude.  The implementation pads the degree to the next power of two using
the maximum representable magnitude.  A balanced winner tournament retains the
direct loser encountered by each winner.  The first minimum is the final
winner; the second minimum is a balanced minimum over the winner's direct
losers.  Stable left preference is expressed with a less-than-or-equal
comparison.  Duplicate values are retained by construction.

For $P$ padded inputs, the structure uses approximately
$P+\log_2P-2$ comparisons: $P-1$ in the tournament and
$\log_2P-1$ along the loser path.  This is a structural argument; FPGA timing
must still be established by static timing analysis.

\subsection{Scaling policies}
The Valls policy constrains
\begin{equation}
 \alpha=2^{-a}+2^{-b},
\end{equation}
so each minimum is scaled once using two constant shifts and one addition.
The result is saturated to the edge magnitude width.

The iteration-ramp policy uses
\begin{equation}
 \alpha_t x = x-(x\gg t).
\end{equation}
Only a bounded set of constant-shift alternatives is elaborated, followed by a
small selector.  No floating-point arithmetic, exponentiation, general
divider, or multiplier is inferred.  Scaling the two shared minima before
routing avoids duplicating the operation per edge.

The CNU output is one broadcast minimum pair plus per-edge sign and selector
metadata.  This interface represents the routing structure directly instead
of copying the same minimum pair into every edge bundle.

\section{Variable-Node Units}
\subsection{Shared variable datapath}
Each CNU contribution selects one member of its minimum pair and converts the
result from sign--magnitude to a sign-extended two's-complement value.  A
balanced adder tree computes the marginal once.  Each outgoing extrinsic
message is then
\begin{equation}
 R_{i,j}=M_j-\sigma_{i,j}.
\end{equation}
This total-minus-own construction avoids one independent adder tree per edge.
The hard decision is taken from the full working-width total, after which the
stored marginal and edge outputs are saturated independently.

The base VNU accepts a signed prior port.  The Valls preset ties this prior to
the fixed representation of $+1$; retaining the general port lets the same
datapath support nonuniform priors and the memory-adjusted VNU.

\subsection{Relay memory}
The deterministic Relay bias accepts $\Lambda_j(0)$, the previous saturated
marginal, and an unsigned fixed-point $\beta_j$.  Both products in
Eq.~\eqref{eq:beta} are evaluated at expanded precision, shifted arithmetically
by the fractional width, combined, and saturated once into the accumulator
format.  The resulting bias enters the shared variable datapath.

The registered Relay wrapper retains the latest saturated marginal.  An
\texttt{initialize} event substitutes the prior for the previous marginal,
which makes the first update independent of stale state.  Changing $\beta_j$
without initialization starts a new memory policy while preserving the relayed
posterior.  A pseudorandom generator is deliberately absent: deterministic
coefficients enable exact reproduction, and an RNG can later be tested as an
independent component.

\section{Cycle Organization}
The registered CNU wrapper captures one check update in one cycle.  The
registered VNU wrapper consumes those results and captures one variable update
in the following cycle.  A complete flooding-schedule iteration therefore has
a two-cycle node latency, consistent with the organizations in
Refs.~\cite{valls2021,maurer2025}.

The present \texttt{MinSumIteration2x2} module elaborates
\begin{equation}
 H=\begin{bmatrix}1&1\\1&1\end{bmatrix}
\end{equation}
as a minimal wiring and latency example.  It is a single-iteration datapath,
not a decoder controller.  The input and output edge arrays are check-major,
while the two VNUs gather one edge from each check.

\section{Systolic Gauss--Jordan Mesh}
The \texttt{GaussJordan} subtree implements independent one-bit column and
diagonal elements without BP message or configuration types.  A future
controller can therefore schedule the mesh after BP without coupling their
datapaths.  The fixed opcode encoding is \texttt{PASS}=00,
\texttt{SWAP}=01, \texttt{ADD}=10, and \texttt{LOCK}=11.

For stored bit $r$ and input $d$, the column element passes $d$, swaps or locks
by emitting $r$ and storing $d$, and adds by emitting $r\oplus d$.  The
diagonal element gives \texttt{reduce\_sig\_i}\(\land r\) priority and emits
\texttt{SWAP}; otherwise a zero passes, the first one locks $r$, and later ones
emit \texttt{ADD}.  State, data, and opcode registers use active-high
synchronous reset and hold when disabled.  The reduction output is instead
combinational, reset-masked, and not enable-gated.

The Scala-parameterized mesh contains $n$ rows and $n+\ell$ columns.  Diagonal
cells launch horizontal opcodes, column cells apply them, and the inactive
lower-left triangle is tied off.  Data moves vertically while the reduction
signal advances between diagonal cells through an enabled pipeline of fixed
elaboration-time depth.  Packed $A$ and lifted-$B$ state exports use row-major
order.  The mesh configuration is immutable Scala data, so it adds no runtime
configuration network and remains independent of BP message types.

\section{Implementation Structure}
Table~\ref{tab:files} maps the implementation to its research purpose.

\begin{table}[t]
\caption{Source organization.}
\label{tab:files}
\centering
\begin{tabular}{@{}ll@{}}
\toprule
File & Responsibility \\
\midrule
\texttt{Config.scala} & Validated immutable parameters \\
\texttt{graph/TannerGraph.scala} & Normalized static graph and edge IDs \\
\texttt{TannerNodeGraphs.scala} & Rooted CNU and VNU port views \\
\texttt{Protocol.scala} & Hardware message bundles \\
\texttt{Arithmetic.scala} & Conversion, scaling, sums, saturation \\
\texttt{TwoMin.scala} & Direct-loser tournament \\
\texttt{CheckNode.scala} & Combinational and registered CNU \\
\texttt{VariableNode.scala} & Base and Relay VNUs \\
\texttt{Iteration.scala} & Fixed $2\times2$ example \\
\texttt{GaussJordan/} & Binary systolic elements and mesh \\
\texttt{Generate.scala} & MLIR and SystemVerilog emission \\
\texttt{Reference.scala} & Independent integer oracle \\
\bottomrule
\end{tabular}
\end{table}

The project pins Scala 2.13.18 and Chisel 7.7.0.  Mill is the sole build
interface.  Chisel elaboration produces a FIRRTL-dialect MLIR artifact, and the
bundled CIRCT firtool lowers the same generator to SystemVerilog.  Selected
artifacts should be archived with the generator configuration, Git revision,
and tool versions rather than treating a mutable build directory as a result.

\section{Verification Methodology}
The verification hierarchy is intentionally redundant at semantic boundaries
but avoids duplicate hardware implementations.

\subsection{Independent reference model}
The Scala oracle represents messages with ordinary integers and booleans.  It
does not call Chisel arithmetic helpers.  This separation prevents a shared
implementation mistake from making the DUT and oracle agree accidentally.

\subsection{Unit tests}
The current tests cover:
\begin{itemize}
  \item exhaustive two-minimum behavior for four three-bit inputs;
  \item exhaustive CNU signs, syndromes, magnitudes, zero, and ties for a small node;
  \item directed Valls and iteration-ramp scaling vectors;
  \item fixed-seed randomized VNU comparison against the oracle;
  \item deterministic Relay bias values including $\beta<1$, $\beta=1$, and $\beta>1$;
  \item registered Relay marginal feedback;
  \item pure Scala graph normalization, transpose construction, stable edge
        numbering, and malformed-adjacency rejection;
  \item exact Steane CNU/VNU node views and a two-way bijection with all twelve
        global edges;
  \item exhaustive Gauss--Jordan transitions, reset, enable hold, and
        reduction forwarding;
  \item cycle-accurate mesh topology, reduction delay, and packed-state
        ordering for small parameter sets; and
  \item cycle-accurate comparison of the fixed $2\times2$ graph.
\end{itemize}
ChiselSim uses Verilator for these behavioral tests.  The emitted
SystemVerilog is additionally checked with standalone Verilator lint.  The
single-file CIRCT output contains multiple module declarations, and fixed-point
products intentionally discard fractional bits.  Filename, unused-bit, and
intentional empty boundary-pin warning classes are disabled explicitly; other
warning classes remain fatal.

\subsection{Claims not established by simulation}
Behavioral simulation establishes values and cycle latency.  It does not
establish maximum clock frequency, critical-path delay, lookup-table count,
register count, routing congestion, or power.  Those quantities require a
documented synthesis and implementation flow for a named FPGA and tool
version.  No timing or resource claim is made in the current manuscript.

\section{Reproduction}
From the repository root, the current milestone is reproduced by
\begin{verbatim}
./scripts/test.sh
./scripts/emit.sh static-steane
./scripts/verify-rtl.sh build/generated/static-steane/StaticTannerDatapath.sv StaticTannerDatapath
./examples/GaussJordan/run.sh
cd examples/GaussJordan/TrapezoidMesh
./run.sh
\end{verbatim}
Emission writes inspectable MLIR and SystemVerilog into the named build
directory.  Focused tests may be run with Mill's \texttt{testOnly} target.
Build directories are excluded from version control.

\section{Planned Extensions}
The next implementation stages are deliberately outside this milestone:
\begin{enumerate}
  \item add the Relay controller for iteration count, legs, convergence
        response, and correction output;
  \item reproduce the paper-specific reduced-precision Relay multiplier as a
        compile-time alternative to the exact fixed-point baseline;
  \item add optional, independently verified coefficient generation;
  \item stream multiple syndrome shots while explicitly banking or replicating
        decoder state; and
  \item report post-synthesis and post-route timing, resources, and power.
\end{enumerate}

\section{Threats to Validity}
The static Steane graph validates composition but is not representative of
qLDPC routing pressure.  Small exhaustive widths expose semantic corner cases but do
not prove all parameter combinations.  The exact Relay fixed-point baseline
does not yet reproduce every approximation in the published FPGA datapath.
Finally, CIRCT transformations may change the textual RTL without changing its
behavior, so semantic tests are treated as primary and textual checks are kept
minimal.

\section{Conclusion}
The current \texttt{chipsLDPC} milestone demonstrates a compact Chisel
decomposition of two related qLDPC min-sum architectures.  Shared sign,
minimum, accumulation, extrinsic, and saturation logic is separated from
static scaling and memory policies, while binary Gauss--Jordan elements remain
an independent post-processing component.  Numeric semantics and cycle
behavior are executable rather than implicit.  This provides a controlled base
for graph generation and FPGA experiments without committing the repository to
a large framework prematurely.

\begin{thebibliography}{9}
\bibitem{valls2021}
J.~Valls, F.~Garcia-Herrero, N.~Raveendran, and B.~Vasi\'{c},
``Syndrome-based min-sum versus OSD-0 decoders: FPGA implementation and
analysis for quantum LDPC codes,'' \emph{IEEE Access}, vol.~9, 2021,
doi: 10.1109/ACCESS.2021.3118544.

\bibitem{maurer2025}
T.~Maurer et al., ``Real-time decoding of the gross code memory with FPGAs,''
arXiv:2510.21600, 2025.

\bibitem{steane1996}
A.~M. Steane, ``Error correcting codes in quantum theory,''
\emph{Phys. Rev. Lett.}, vol.~77, pp.~793--797, 1996,
doi: 10.1103/PhysRevLett.77.793.

\bibitem{chiselbook}
M.~Schoeberl, \emph{Digital Design with Chisel}, 6th ed., 2025.

\bibitem{chisel}
CHIPS Alliance, ``Chisel documentation,''
\url{https://www.chisel-lang.org/docs}, accessed Aug.~31, 2026.
\end{thebibliography}

\end{document}
